% Section 1.4: Indirect Detection via Gamma-Rays (The Formalism)
% Structure: Annihilation/decay physics -> Spectral signatures -> Density profiles -> J/D-factor -> Targets -> Multi-messenger -> Status

\section{Indirect Detection via Gamma-Rays}
\label{sec:1.4}

Section~\ref{sec:1.3} established that indirect detection via gamma-rays directly probes the annihilation cross-section --- the quantity that governs the thermal relic abundance and the one quantity that neither underground laboratories nor colliders can access.
This section develops the underlying formalism.
We follow the signal from particle physics (what happens when two dark matter particles meet) through astrophysics (where the dark matter sits in the sky) to yield the master equations that predict the gamma-ray flux from any given target.
These equations, and the assumptions that enter them, underpin every analysis presented in this thesis.

The development begins with the microphysics of annihilation and decay and the gamma-ray spectra they produce (\S\ref{sec:ann_decay}--\S\ref{sec:spectral_features}), then introduces the dark matter density profiles that determine how signal strength varies across the sky (\S\ref{sec:density_profiles}).
These ingredients combine into the $J$-factor and $D$-factor formalism --- a factorization of the predicted flux into a particle physics term and an astrophysical term (\S\ref{sec:jfactor}) --- before we survey the observational targets that define the thesis programme (\S\ref{sec:targets}--\S\ref{sec:status}).

\subsection{Annihilation and Decay Physics}
\label{sec:ann_decay}

Dark matter particles in present-day galactic halos are cold --- their typical velocities are $v \sim 10^{-3}\, c$ --- and annihilate or decay essentially at rest \cite{Cirelli:2024ssz}.
This kinematic fact has two immediate consequences.
First, the total center-of-mass energy available in an annihilation event is simply $\sqrt{s} = 2 m_\chi$, so the most energetic photon that can be produced carries energy $E_\gamma = m_\chi$ \cite{Cirelli:2024ssz}.
For decay, the available energy is instead $m_\chi$, and momentum conservation among at least two daughter particles limits the maximum photon energy to $E_\gamma = m_\chi / 2$ \cite{Cirelli:2024ssz,Pinetti:2021jjs}.
These hard kinematic endpoints are one of the most distinctive features of a dark matter signal: no astrophysical process produces a gamma-ray spectrum with such a sharp cutoff tied to a single mass scale.

The second consequence is more profound and concerns the spatial morphology of the expected signal.
Annihilation is a two-body process --- two dark matter particles must find each other --- and its rate per unit volume is therefore proportional to the square of the local dark matter density:

\begin{equation}
	\frac{d\Gamma_\text{ann}}{dV} = \frac{\langle \sigma v_\text{rel} \rangle}{2} \frac{\rho^2}{m_\chi^2} \,,
	\label{eq:ann_rate}
\end{equation}

where $\langle \sigma v_\text{rel} \rangle$ is the thermally averaged annihilation cross-section and the factor of $1/2$ avoids double-counting for self-conjugate particles \cite{Hooper:2024,Cirelli:2024ssz}.
(For Dirac dark matter, where particle and antiparticle are distinct, the factor becomes $1/4$ but the thermal cross-section required for the correct relic abundance is correspondingly larger, leaving the observable rate unchanged \cite{Hooper:2024}.) Decay, by contrast, is a single-particle process with a rate proportional to the first power of the density:

\begin{equation}
	\frac{d\Gamma_\text{dec}}{dV} = \frac{\rho}{m_\chi \tau} \,,
	\label{eq:dec_rate}
\end{equation}

where $\tau$ is the dark matter lifetime \cite{Hooper:2024,Pinetti:2021jjs}.
Current gamma-ray observations constrain $\tau \gtrsim 10^{28}$ s \cite{Cirelli:2024ssz}, far exceeding the age of the universe --- making decay an intrinsically rare process whose detection requires enormous integrated exposures.
The $\rho^2$ versus $\rho$ distinction has direct observational consequences: annihilation signals are sharply concentrated in the densest regions of a halo --- the inner kiloparsec of the Galactic Center, or the cores of dwarf galaxies --- while decay signals are spread diffusely across the entire halo volume, producing a much flatter angular profile \cite{Cirelli:2024ssz}.
This morphological difference provides a powerful diagnostic: an excess that tracks the $\rho^2$ profile of the inner Galaxy but vanishes at high latitudes points toward annihilation, while a persistent, broadly distributed signal favours decay.

The stable particles emerging from an annihilation or decay event --- photons, neutrinos, electrons, positrons, protons, antiprotons --- are determined by the primary channel into which the dark matter annihilates.
In a model-independent analysis, one parametrizes the annihilation as $\chi\chi \to f\bar{f}$ for a specific Standard Model final state $f$ (quarks, leptons, or gauge bosons), and the total gamma-ray spectrum is given by a sum over channels weighted by branching ratios \cite{Cirelli:2024ssz}:

\begin{equation}
	\frac{dN_\gamma}{dE} = \sum_f \text{BR}_f \, \frac{dN_\gamma^{(f)}}{dE} \,.
	\label{eq:spectrum_sum}
\end{equation}

In practice, limits are derived for one channel at a time --- setting $\text{BR}_f = 1$ for a chosen final state --- so each "benchmark channel" maps to a distinctive spectral shape.
The spectra $dN_\gamma^{(f)}/dE$ for each channel are computed using Monte Carlo event generators such as Pythia, with the results tabulated in standard libraries such as PPPC 4 DM ID \cite{Cirelli:2010xx,Hooper:2024}.
These spectra exhibit qualitatively different shapes depending on $f$.
In hadronic channels such as $b\bar{b}$ or $t\bar{t}$, the quarks hadronize into jets containing copious neutral pions, which decay almost instantly via $\pi^0 \to \gamma\gamma$.
The resulting gamma-ray spectrum is broad and soft, peaking at $E_\gamma \sim m_\chi / 20$ in $E^2 \, dN/dE$ units \cite{Hooper:2024,Cirelli:2024ssz}.
The $b\bar{b}$ channel has become the standard benchmark in Fermi-LAT analyses because it produces a generic spectrum representative of many WIMP models.

Leptonic channels behave differently.
The $\tau^+\tau^-$ final state undergoes hadronic decays that produce fewer but more energetic pions, yielding a harder spectrum peaking at $E_\gamma \sim m_\chi / 3$ \cite{Hooper:2024}.
Lighter leptons ($e^+e^-$, $\mu^+\mu^-$) do not produce pions at tree level --- their gamma-ray spectrum is dominated by final-state radiation and carries far less total flux, concentrated near the kinematic endpoint \cite{Cirelli:2024ssz}.
Gauge boson channels ($W^+W^-$, $ZZ$) produce a combination of both behaviours: the hadronic decays of the $W$ and $Z$ generate a broad low-energy bump similar to $b\bar{b}$, while the charged $W^\pm$ bosons also radiate hard photons directly, adding a hard photon tail near $E_\gamma \sim m_\chi$ \cite{Cirelli:2024ssz}.

The underlying pattern is simple: the more hadrons produced in the cascade, the softer the spectrum, because the energy is shared among a large multiplicity of pions.
Channels that produce fewer, more energetic particles yield harder spectra with sharper features.
In practice, the distinction between the $b\bar{b}$ and $\tau^+\tau^-$ benchmarks brackets the range of spectral shapes relevant for most thermal relic models \cite{Hooper:2024}.
These qualitative differences are illustrated in Figure~\ref{fig:spectra_per_channel}.

\begin{figure}[t]
	\centering
	\includegraphics[width=0.48\textwidth]{figures/Daylan2014_spectra_per_channel.pdf}
	\caption{
	Prompt gamma-ray spectrum $dN_\gamma/dE_\gamma$ produced per dark matter annihilation for a 30~GeV WIMP, shown for a variety of Standard Model final states.
		Hadronic channels ($b\bar{b}$, $c\bar{c}$) produce soft, broad spectra peaked at $E_\gamma \sim m_\chi/20$, while leptonic channels ($\tau^+\tau^-$, $\mu^+\mu^-$) yield harder spectra concentrated near the kinematic endpoint.
		Figure from Daylan et al.~\cite{Daylan:2014rsa}.
	}
	\label{fig:spectra_per_channel}
\end{figure}

\subsection{Spectral Features and Signatures}
\label{sec:spectral_features}

The continuum spectra described above are broad, featureless apart from a cutoff at $m_\chi$, and not unique to dark matter --- astrophysical gamma-ray sources produce similar power-law-like spectra through conventional particle acceleration \cite{Cirelli:2024ssz}.
The discriminating power of gamma-ray searches therefore lies in spectral features that have no plausible astrophysical counterpart.
Unlike charged cosmic rays, whose spectral features are smeared by propagation through galactic magnetic fields, sharp structures in the gamma-ray spectrum propagate intact from source to detector and can be searched for directly \cite{Cirelli:2024ssz}.

\textbf{Monochromatic lines.}
The cleanest such feature is a gamma-ray line: a monochromatic photon emitted at a fixed energy $E_\gamma = m_\chi$ from the process $\chi\chi \to \gamma\gamma$ \cite{Cirelli:2024ssz}.
Because dark matter is electrically neutral, this annihilation cannot proceed at tree level and must be mediated by a loop of charged particles, suppressing the cross-section by a loop factor relative to direct annihilation into Standard Model final states \cite{Cirelli:2024ssz,Pinetti:2021jjs}.
The expected flux is correspondingly small.
Yet no known astrophysical mechanism produces a sharp line superimposed on smooth backgrounds, which is why the monochromatic signal has long been regarded as a ``smoking gun'' for dark matter \cite{Cirelli:2024ssz}.
A related process, $\chi\chi \to \gamma Z$ or $\gamma h$, produces a secondary line displaced to lower energy, $E_\gamma = m_\chi(1 - M_{Z,h}^2/4m_\chi^2)$, providing an additional handle for identification \cite{Cirelli:2024ssz}.
An interesting exception to the generic loop suppression arises in Minimal Dark Matter models, where the non-perturbative Sommerfeld enhancement boosts the loop-induced cross-section to a level comparable with the continuum contribution for a range of dark matter masses \cite{Cirelli:2024ssz}.

\textbf{Virtual internal bremsstrahlung (VIB).}
A second sharp feature arises when a photon is radiated from a virtual charged particle mediating the annihilation \cite{Bringmann:2007nk}.
Because this emission brings the virtual propagator closer to its mass shell --- particularly when the mediator mass is close to $m_\chi$ --- the process preferentially produces hard photons near the kinematic limit \cite{Cirelli:2024ssz}.
The resulting VIB spectrum has a characteristic saw-tooth shape: rising steeply toward a sharp cutoff at $E_\gamma \lesssim m_\chi$ \cite{Bringmann:2007nk,Cirelli:2024ssz}.
Beyond its distinctive spectral signature, VIB has an important kinematic consequence.
As a three-body final state, it lifts the helicity suppression that affects two-body annihilation of Majorana fermions into light fermions: the emission of the additional photon removes the strict angular momentum constraints on the two-body process, allowing the annihilation to proceed via unsuppressed $s$-wave \cite{Bringmann:2007nk,Cirelli:2024ssz}.
In models where the two-body rate is helicity-suppressed, the VIB cross-section can actually exceed the tree-level two-body rate, making the three-body radiative process the phenomenologically dominant channel \cite{Cirelli:2024ssz}.
These spectral features are illustrated schematically in Figure~\ref{fig:spectral_features}.

\begin{figure}[t]
	\centering
	\includegraphics[width=0.42\textwidth]{figures/IndirectDetectionSpectrumCartoonNew.pdf}
	\caption{Notable sharp gamma-ray spectral features from dark matter annihilation.
		From left to right: the monochromatic line from $\chi\chi \to \gamma\gamma$, the ``saw-tooth'' shape of virtual internal bremsstrahlung (VIB), the ``top-hat'' (box-shaped) spectrum from cascade decays $\chi\chi \to XX \to \gamma\gamma\gamma\gamma$, and the single saw-tooth from $\chi\chi \to XX \to \gamma\nu\gamma\nu$.
		Unlike charged cosmic-ray features, these sharp structures propagate intact to the detector and can be searched for directly.
		Figure from Cirelli et al.~\cite{Cirelli:2024ssz}.
	}
	\label{fig:spectral_features}
\end{figure}

\textbf{Electroweak corrections.}
For TeV-scale dark matter ($m_\chi \gg M_W$), the radiation of $W$ and $Z$ bosons from the annihilation products becomes significant \cite{Ciafaloni:2010ti,Cirelli:2024ssz}.
Although this is formally a higher-order process suppressed by the weak coupling, large logarithmic enhancements of the form $\ln(m_\chi / M_W)$ compensate the suppression at high masses \cite{Cirelli:2024ssz}.
The consequences are twofold.
First, the sharp peaks of monochromatic lines are smeared, and a new continuum of lower-energy gamma-rays ($E_\gamma \ll m_\chi$) appears from the hadronic and leptonic decays of the radiated gauge bosons \cite{Cirelli:2024ssz}.
Second, purely leptonic channels ($\chi\chi \to e^+e^-$, $\mu^+\mu^-$) and even the $\gamma\gamma$ channel acquire an antiproton yield that would be entirely absent at tree level, altering the multi-messenger predictions \cite{Cirelli:2024ssz}.
These electroweak corrections were incorporated into the standard spectral libraries only after 2013 \cite{Cirelli:2010xx,Cirelli:2024ssz}, and neglecting them leads to qualitatively wrong predictions for heavy dark matter candidates.

\textbf{Prompt versus secondary emission.}
The gamma-ray production mechanisms discussed so far --- pion decay, final-state radiation, lines, VIB --- all occur directly in the annihilation or decay event and are collectively termed \textit{prompt emission} \cite{Cirelli:2024ssz,Pinetti:2021jjs}.
Prompt photons propagate freely through the Galaxy and arrive at the detector with their spectral shape intact, because the Galaxy is effectively transparent to gamma-rays below a few hundred TeV \cite{Hooper:2024}.

A qualitatively different contribution arises when the energetic electrons and positrons produced by dark matter propagate through the interstellar medium.
These leptons upscatter ambient low-energy photons to gamma-ray energies via \textit{inverse Compton scattering} (ICS) \cite{Cirelli:2024ssz,Pinetti:2021jjs}.
The target photon field --- the interstellar radiation field (ISRF) --- consists of three components: the cosmic microwave background, optical starlight, and dust-reprocessed infrared light \cite{Cirelli:2024ssz,Pinetti:2021jjs}.
To fix orders of magnitude, a 1 TeV electron upscatters a CMB photon to $\sim 1.5$ GeV, or a UV starlight photon to $\sim 150$ GeV \cite{Cirelli:2024ssz}.
The propagating $e^\pm$ also emit \textit{bremsstrahlung} radiation when deflected by the electric fields of interstellar gas nuclei, a process that dominates their energy losses in the 1--10 GeV range near the Galactic plane where gas densities are highest \cite{Cirelli:2024ssz}.
These processes produce \textit{secondary} gamma-ray emission that is spatially extended and spectrally softer than the prompt component, and whose morphology depends on the magnetic field structure, gas distribution, and radiation fields of the host galaxy \cite{Cirelli:2024ssz,Pinetti:2021jjs}.
Secondary emission is sometimes called ``doubly indirect'' because it depends not only on the dark matter distribution but also on the astrophysical environment that reprocesses the electrons \cite{Cirelli:2024ssz}.
Despite this additional layer of modelling uncertainty, secondary emission plays a central role in the interpretation of the Galactic Center excess discussed in Chapter~\ref{ch:4}.
The very identification of the excess relies on subtracting the diffuse gamma-ray background dominated by cosmic-ray ICS and bremsstrahlung in the dense inner Galaxy \cite{Cirelli:2024ssz,Hooper:2024}.
Indeed, ICS constraints from peaked density profiles have been used to exclude leptophilic dark matter interpretations of the signal \cite{Cirelli:2024ssz}.

\subsection{How Dark Matter is Distributed: Density Profiles}
\label{sec:density_profiles}

The annihilation and decay rates derived in Section~\ref{sec:ann_decay} depend directly on the dark matter density $\rho(r)$ --- and for annihilation, on its square.
The predicted gamma-ray flux from any astrophysical target therefore hinges on the assumed density profile.
The same measurement that constrains the particle physics simultaneously probes the dark matter distribution, but the result is only as reliable as the profile assumed --- a tension that runs through every chapter of this thesis.

The standard profiles used in the field come from cosmological $N$-body simulations --- numerical experiments that follow the gravitational collapse of billions of dark matter particles from primordial density perturbations to the present-day cosmic web \cite{Cirelli:2024ssz}.
These simulations show that the spherically averaged density $\rho(r)$ of collapsed halos follows a universal logarithmic slope $s(r) \equiv d\ln\rho / d\ln r$ that transitions smoothly from $s \approx -1$ in the inner region to $s \approx -3$ in the outskirts \cite{Cirelli:2024ssz}.
The results can be captured by several fitting functions, each with a characteristic density $\rho_s$ and a scale radius $r_s$.

\textbf{NFW (Navarro--Frenk--White).}
The most widely adopted benchmark has a density that diverges as $\rho \propto r^{-1}$ toward the center (a ``cusp'') and falls off as $\rho \propto r^{-3}$ in the outer halo \cite{Navarro:1996gj}:

\[ \rho_\text{NFW}(r) = \frac{\rho_s}{\dfrac{r}{r_s}\left(1 + \dfrac{r}{r_s}\right)^2} \,, \]

A common generalization replaces the fixed inner slope with a free parameter $\gamma$, giving the ``generalized NFW'' (gNFW) profile $\rho \propto (r/r_s)^{-\gamma}(1 + r/r_s)^{\gamma - 3}$ \cite{Cirelli:2024ssz}.
The standard NFW corresponds to $\gamma = 1$; steeper values ($\gamma > 1$) describe adiabatically contracted halos, while $\gamma < 1$ accommodates shallower cusps.
This freedom is directly relevant for the Galactic Center excess analysis in Chapter~\ref{ch:4}, where the morphology of the signal constrains $\gamma$.

\textbf{Einasto.}
Originally proposed by Einasto (1965) and later identified as a fitting function for simulated halos \cite{Cirelli:2024ssz}, this profile has a logarithmic slope that varies continuously with radius, without converging to a fixed power law at the center.
It is controlled by a shape parameter $\alpha_\text{Ein}$:

\[ \rho_\text{Ein}(r) = \rho_s \exp\left\{ -\frac{2}{\alpha_\text{Ein}} \left[\left(\frac{r}{r_s}\right)^{\alpha_\text{Ein}} - 1\right]\right\} \,.
\]

High-resolution simulations increasingly favour Einasto over NFW as a fitting function, with $\alpha_\text{Ein} \approx 0.17$ providing a good average fit across a range of halo masses \cite{Cirelli:2024ssz}.
At the specific radii probed by indirect detection, however, the two profiles yield similar predictions, and both are used as standard benchmarks.

\textbf{Burkert (cored).}
A profile with a constant-density core at small radii, motivated by observed rotation curves of dwarf and low-surface-brightness galaxies rather than $N$-body simulations \cite{Burkert:1995yz,Bullock:2017xww}:

\[ \rho_\text{Bur}(r) = \frac{\rho_s}{\left(1 + \dfrac{r}{r_s}\right)\left(1 + \left(\dfrac{r}{r_s}\right)^2\right)} \,.
\]

These three profiles are compared in Figure~\ref{fig:dm_profiles}, which shows both the radial density dependence and the fitted parameters for the Milky Way.
At large radii (around $r_\odot$) all profiles agree, but they diverge by orders of magnitude in the inner kiloparsec --- the region most relevant for indirect detection from the Galactic Center.

\begin{figure}[t]
	\centering
	\includegraphics[width=\textwidth]{figures/DMprofiles.pdf}
	\caption{Dark matter density profiles and their Milky Way parameters. \textit{Left}: NFW (red), Einasto (blue), Burkert (green), and Isothermal (purple) profiles as a function of galactocentric distance.
		The profiles agree at $r \sim r_\odot$ but diverge by orders of magnitude in the inner kiloparsec.
		\textit{Right}: the corresponding scale radii $r_s$ and characteristic densities $\rho_s$ in the parametrizations of each profile.
		The precision shown ($r_s$ to 2 and $\rho_s$ to 3 significant digits) is sufficient for most computations; more precise inputs are needed to reproduce $J$-factors for small angular regions around the Galactic Center.
		Figure from Cirelli et al.~\cite{Cirelli:2024ssz}.
	}
	\label{fig:dm_profiles}
\end{figure}

When fit to the Milky Way, the parameters of these profiles are fixed by requiring consistency with the total Galactic mass ($\sim 4.7 \times 10^{11}\, M_\odot$) and the local dark matter density at the Sun's position ($r_\odot \simeq 8.12$ kpc) \cite{Cirelli:2024ssz,Pinetti:2021jjs}.
The most recent determinations place the local density at $\rho_\odot \approx 0.40 \pm 0.10$ GeV/cm$^3$ \cite{Cirelli:2024ssz}, though many analyses continue to use the historical benchmark of 0.3 GeV/cm$^3$ for consistency with existing limits.
Characteristic Milky Way parameters are $r_s \approx 14.5$ kpc and $\rho_s \approx 0.57$ GeV/cm$^3$ for NFW, and $r_s \approx 13.6$ kpc with $\rho_s \approx 0.15$ GeV/cm$^3$ for Einasto \cite{Cirelli:2024ssz}.

At the solar radius, all three profiles agree by construction.
They differ sharply in the inner kiloparsec, where no direct observational data constrain the dark matter profile and the resolution of current simulations reaches its limit (roughly a fraction of a kiloparsec for a Milky Way--sized halo) \cite{Cirelli:2024ssz}.
The NFW and Einasto profiles predict densities at $r \sim 100$ pc that exceed the Burkert core prediction by one to two orders of magnitude \cite{Cirelli:2024ssz}.
This discrepancy propagates directly into the $J$-factor: the predicted gamma-ray flux from the Galactic Center can vary by orders of magnitude depending on whether one assumes a cusp or a core.

The disagreement between cuspy simulations and cored observations is the \textbf{cusp--core problem} \cite{Cirelli:2024ssz,Bullock:2017xww}.
Dark matter--only $N$-body simulations consistently produce cusps with inner slopes $\gamma \approx 1$--$1.5$, while rotation curve measurements --- particularly in dwarf galaxies --- often favour constant-density cores ($\gamma \approx 0$--$0.5$) \cite{Bullock:2017xww}.
Several resolutions have been proposed.
From the particle physics side, self-interacting dark matter (with $\sigma/m \approx 0.5$--$10$ cm$^2$/g) transfers heat to halo centers and produces cores of the observed size \cite{Cirelli:2024ssz,Bullock:2017xww}, while ultralight ``fuzzy'' dark matter generates solitonic cores through quantum pressure on scales comparable to the de Broglie wavelength \cite{Cirelli:2024ssz}.
From the astrophysical side, modern hydrodynamical simulations show that repeated bursts of supernova-driven outflows can rapidly alter the central gravitational potential, pushing dark matter outward and flattening a primordial cusp into a core \cite{Cirelli:2024ssz,Bullock:2017xww}.

The effectiveness of this baryonic feedback depends on the stellar-to-halo mass ratio $M_\star / M_\text{vir}$ in a non-monotonic way \cite{Cirelli:2024ssz,Bullock:2017xww}.
At very low ratios ($M_\star / M_\text{vir} \lesssim 10^{-4}$, corresponding to ultrafaint dwarfs with $M_\star \lesssim 10^6\, M_\odot$), too few stars form to generate the feedback energy needed to modify the cusp, and the profile remains NFW-like.
Core formation peaks at intermediate ratios ($M_\star / M_\text{vir} \approx 5 \times 10^{-3}$, the regime of bright dwarfs with $M_\star \sim 10^{8}$--$10^{9}\, M_\odot$), where supernovae are frequent enough to redistribute the dark matter without the stellar mass dominating the potential.
At high ratios ($M_\star / M_\text{vir} \gtrsim 0.01$), the concentration of baryons at the galaxy's center deepens the gravitational well and pulls dark matter inward through adiabatic contraction, recreating a cusp.
For the Milky Way ($M_\star / M_\text{vir} \approx 0.08$), this puts the expected inner profile firmly in the cuspy regime \cite{Cirelli:2024ssz} --- but the argument rests on the assumption that feedback and contraction have been correctly modelled, which remains uncertain.

\textbf{Substructure and boost factors.}
Real dark matter halos are not smooth.
$N$-body simulations reveal a rich hierarchy of gravitationally bound subhalos embedded within the main halo, with a mass function $dn/dM \propto M^{-2}$ extending from cluster masses down to the resolution limit \cite{Cirelli:2024ssz,Springel:2008cc}.
Because annihilation scales as $\rho^2$, the clumpy density field produces a higher total signal than a smooth halo of the same average density: $\langle \rho^2 \rangle > \langle \rho \rangle^2$.
This enhancement is parametrized by a multiplicative \textbf{boost factor} $B \equiv \langle \rho^2 \rangle / \langle \rho \rangle^2$ \cite{Cirelli:2024ssz}.

Early estimates of the boost factor ranged from modest ($B \sim 10$) to extreme ($B \sim 10^9$), but modern analyses find much smaller values.
The survival of subhalos is governed by tidal stripping from the host halo's gravitational field, gravitational ``disk shocking'' as subhalos cross the dense baryonic disk, and encounters with individual stars --- all of which preferentially destroy subhalos in the inner Galaxy \cite{Cirelli:2024ssz}.
When these effects are included, the boost factor for the Galaxy as a whole is $B \lesssim \mathcal{O}(10)$--$\mathcal{O}(30)$ \cite{Cirelli:2024ssz}.
Toward the Galactic Center, where subhalos are most heavily depleted by tidal forces, the boost is expected to be negligible \cite{Cirelli:2024ssz}.
The outskirts of the halo, by contrast, retain a richer substructure population, and the individual subhalos that survive there are the targets of the searches described in Chapters~\ref{ch:5} and~\ref{ch:6}.

The subhalo hierarchy does not extend to arbitrarily small masses.
It is truncated at a minimum mass $M_\text{min}$ set by the kinetic decoupling of the dark matter particle from the thermal plasma in the early universe \cite{Cirelli:2024ssz}.
After kinetic decoupling, dark matter particles free-stream over a characteristic length $\lambda_\text{fs}$, erasing density perturbations below the corresponding mass scale $M_\text{min} \simeq (4\pi/3)\, \rho_\text{DM}\, \lambda_\text{fs}^3$ \cite{Cirelli:2024ssz}.
For a 100 GeV WIMP with weak-scale couplings, $M_\text{min} \sim 10^{-6}\, M_\odot$ --- roughly an Earth mass \cite{Cirelli:2024ssz}.
This scale lies far below the resolution of any current simulation and represents a fundamental uncertainty in the predicted total annihilation rate, since the integral over the subhalo mass function is dominated by the smallest masses.

The density profiles and substructure properties introduced here feed directly into the $J$-factor and $D$-factor formalism of Section~\ref{sec:jfactor}, which combines them with the particle physics of Section~\ref{sec:ann_decay} to produce the master equations for gamma-ray indirect detection.

\subsection{The Flux Factorization: $J$-Factor and $D$-Factor}
\label{sec:jfactor}

Combining the particle physics of Section~\ref{sec:ann_decay} with the density profiles of Section~\ref{sec:density_profiles}, we can now write the differential gamma-ray flux from dark matter annihilation observed from a direction $\psi$ relative to the Galactic Center \cite{Cirelli:2024ssz,Hooper:2024}:

\begin{equation}
	\frac{d\Phi_\gamma}{dE\, d\Omega} = \underbrace{\frac{1}{4\pi} \frac{\langle \sigma v_\text{rel} \rangle}{2 m_\chi^2} \frac{dN_\gamma}{dE}}_{\text{Particle Physics}} \times \underbrace{J(\psi)}_{\text{Astrophysics}} \,.
	\label{eq:flux_ann}
\end{equation}

The prefactor $1/(4\pi)$ reflects the isotropic emission of photons over the full sphere, and the factor of $1/2$ avoids double-counting for self-conjugate (Majorana) particles, as discussed in Section~\ref{sec:ann_decay}.
The $J$-factor is the line-of-sight integral of the density squared:

\begin{equation}
	J(\psi) = \int_\text{l.o.s.} \rho^2\bigl(r(s, \psi)\bigr) \, ds \,,
	\label{eq:jfactor}
\end{equation}

where $s$ parametrizes distance along the line of sight and $r$ is the galactocentric radius, related to $s$ and the observation angle $\psi$ by $r = (r_\odot^2 + s^2 - 2 r_\odot s \cos\psi)^{1/2}$, with $r_\odot \simeq 8.33$ kpc the Sun--Galactic Center distance \cite{Cirelli:2024ssz}.
In this dimensionful convention, $J$ carries units of GeV$^2$ cm$^{-5}$; an alternative dimensionless definition, in which $\rho$ and $s$ are normalized to $\rho_\odot$ and $r_\odot$ respectively, is also common in the literature \cite{Cirelli:2024ssz}.

For decaying dark matter, the analogous expression is \cite{Cirelli:2024ssz,Pinetti:2021jjs}:

\begin{equation}
	\frac{d\Phi_\gamma}{dE\, d\Omega} = \frac{1}{4\pi} \frac{\Gamma}{m_\chi} \frac{dN_\gamma}{dE} \times D(\psi) \,,
	\label{eq:flux_dec}
\end{equation}

where $\Gamma = 1/\tau$ is the decay rate.
No symmetry factor appears, since decay is a single-particle process.
The $D$-factor integrates the density linearly:

\begin{equation}
	D(\psi) = \int_\text{l.o.s.} \rho\bigl(r(s, \psi)\bigr) \, ds \,,
	\label{eq:dfactor}
\end{equation}

and carries units of GeV cm$^{-2}$.
Because $D$ depends on $\rho$ rather than $\rho^2$, it is far less sensitive to the assumed density profile than the $J$-factor, and substructure boost factors are entirely absent for decay \cite{Cirelli:2024ssz}.
This robustness makes decay searches less model-dependent than their annihilation counterparts, though at the cost of weaker expected signals.

The power of this factorization is that it cleanly separates the unknowns.
The particle physics factor --- containing $m_\chi$, $\langle \sigma v_\text{rel} \rangle$ (or $\tau$), and the spectral shape $dN_\gamma/dE$ --- encodes the dark matter model we wish to test.
The astrophysical factor --- $J$ or $D$ --- encodes the distribution of dark matter along the line of sight, determined by the density profile of the host halo and constrained independently through kinematic measurements.
A non-detection in a target with a well-determined $J$-factor translates directly into an upper limit on $\langle \sigma v_\text{rel} \rangle$; for decay, an upper limit on $\Gamma$ (equivalently, a lower limit on $\tau$).

When applied to an extended region of the sky, the integrated flux is obtained by integrating $J$ (or $D$) over the solid angle of interest.
For comparing different targets, one typically quotes the $J$-factor integrated over a cone of half-angle $\theta_\star$ centered on the target:

\begin{equation}
	J_\text{int}(\theta_\star) = \int_0^{\theta_\star} J(\psi) \, 2\pi \sin\psi \, d\psi \,.
	\label{eq:jfactor_int}
\end{equation}

For dwarf spheroidal galaxies observed by Fermi-LAT, typical integrated $J$-factors (within $\theta_\star = 0.5^\circ$) range from $\sim 10^{16}$ to $\sim 10^{19.5}$ GeV$^2$ cm$^{-5}$, depending on the system's mass, distance, and the quality of the kinematic data used to reconstruct its density profile \cite{Hooper:2024,Cirelli:2024ssz,Bonnivard:2015xpq}.
These values are orders of magnitude smaller than the Galactic Center $J$-factor, but the absence of astrophysical backgrounds in dwarfs compensates for the reduced signal.

Several caveats affect the simple factorization above.
First, it assumes a \textit{velocity-independent} annihilation cross-section.
If $\langle \sigma v_\text{rel} \rangle$ depends on the dark matter velocity --- as it does for $p$-wave annihilation ($\sigma v \propto v^2$) or Sommerfeld-enhanced models ($\sigma v \propto 1/v$) --- the cross-section can no longer be factored out of the line-of-sight integral \cite{Cirelli:2024ssz}.
In such cases, one must compute an \textit{effective $J$-factor} that folds the full, position-dependent velocity distribution into the integral, and the clean particle-physics/astrophysics separation is lost.

Second, even for velocity-independent cross-sections, the $J$-factor depends quadratically on the density profile, making it extremely sensitive to the uncertainties discussed in Section~\ref{sec:density_profiles}.
For the Galactic Center, the $J$-factor computed with an NFW profile can exceed the Burkert value by orders of magnitude \cite{Cirelli:2024ssz}; even within the NFW family, varying the inner slope from $\gamma = 0.7$ to $\gamma = 1.3$ changes the predicted flux by more than an order of magnitude \cite{Hooper:2024}.
The astrophysical factor is therefore the dominant systematic uncertainty in annihilation searches --- a theme that recurs in every chapter of this thesis.

\subsection{Observational Targets for Gamma-Ray Searches}
\label{sec:targets}

The formalism of the preceding sections produces a clear prescription: to maximise the gamma-ray signal from dark matter, one should observe targets with large $J$-factors (for annihilation) or $D$-factors (for decay) and low astrophysical backgrounds.
No single target satisfies both criteria simultaneously, and the field has consequently developed a multi-target observation programme in which the complementary strengths of different targets compensate each other's weaknesses \cite{Cirelli:2024ssz}.
We now review the principal targets, noting their advantages, limitations, and connection to the analyses developed in this thesis.

\textbf{The Galactic Center.}
The center of the Milky Way offers the largest $J$-factor of any astronomical source, owing to its proximity ($\sim 8$ kpc) and the expected steep rise of the dark matter density at small radii \cite{Hooper:2024,Pinetti:2021jjs}.
For an NFW profile with $\gamma = 1$ and a local density $\rho_\odot = 0.4$ GeV/cm$^3$, the predicted flux from the inner $10^\circ$ exceeds that from any individual dwarf galaxy by more than three orders of magnitude \cite{Hooper:2024}.
If a steeper inner slope ($\gamma \approx 1.3$) is assumed, the flux increases by an additional factor of $\sim 5$; for a shallower slope ($\gamma \approx 0.7$) it decreases by $\sim 3$ \cite{Hooper:2024}.
In either case, the Galactic Center remains the single brightest expected source of dark matter annihilation products on the sky.

The advantage is offset by severe astrophysical contamination.
The inner few degrees of the Galaxy host bright diffuse emission from cosmic-ray protons and electrons scattering off interstellar gas (through pion production and bremsstrahlung) and ambient radiation fields (through inverse Compton scattering), together with numerous point sources --- supernova remnants, pulsars, and the central supermassive black hole \cite{Hooper:2024,Cirelli:2024ssz}.
Models of this diffuse foreground are constructed from gas maps and the locally measured cosmic-ray spectrum, and while they capture the broad features of the observed emission, they cannot account for its detailed spectral and morphological structure \cite{Hooper:2024}.
This imperfect background knowledge is the primary obstacle to extracting a dark matter signal from the Galactic Center.

Despite these difficulties, the Fermi-LAT has detected a statistically significant excess of gamma-rays from the inner Galaxy at energies between $\sim 0.5$ and 5 GeV --- the so-called Galactic Center excess (GCE) \cite{Cirelli:2024ssz,Daylan:2014rsa}.
The spectrum and approximately spherical morphology of this excess are consistent with the predictions for a $\sim 40$--$70$ GeV dark matter particle annihilating into $b\bar{b}$ quarks at a rate close to the thermal relic value \cite{Cirelli:2024ssz,Daylan:2014rsa}.
The dark matter interpretation is tantalising, but an unresolved population of millisecond pulsars in the Galactic bulge provides an equally plausible astrophysical explanation \cite{Cirelli:2024ssz,Pinetti:2021jjs}.
Multiple statistical analyses have investigated the point-source versus diffuse origin of the excess, reaching contradictory conclusions \cite{Cirelli:2024ssz}.
The morphological debate --- whether the excess traces a spherical dark matter halo or the stellar distribution of the Galactic bulge --- remains unresolved \cite{Cirelli:2024ssz}.
This persistent ambiguity motivates the dedicated analysis in Chapter~\ref{ch:4}, which examines the millisecond pulsar hypothesis by constraining the pulsar luminosity function from globular cluster observations.
Figure~\ref{fig:gce} shows the residual gamma-ray maps from the inner Galaxy after subtraction of the standard diffuse model, revealing the spatially extended excess at 0.5--5~GeV, and the spectrum of this excess compared with the prediction for annihilating dark matter.

\begin{figure}[t]
	\centering
	\begin{minipage}[t]{0.48\textwidth}
		\centering
		\includegraphics[width=\textwidth]{figures/Daylan2014_residual_map.pdf}
	\end{minipage}
	\hfill
	\begin{minipage}[t]{0.48\textwidth}
		\centering
		\includegraphics[width=\textwidth]{figures/spectrum-CCWR1pt18.pdf}
	\end{minipage}
	\caption{\textit{Left}: Residual gamma-ray intensity maps (in Galactic coordinates) after subtracting the best-fit Galactic diffuse model, Fermi bubbles, and isotropic templates.
		The dark-matter-like excess is clearly visible around the Galactic Center at energies between $\sim$0.5 and 5~GeV.
		\textit{Right}: Spectrum of the excess for a generalized NFW template with inner slope $\gamma = 1.18$, compared with the prediction for a 43~GeV dark matter particle annihilating to $b\bar{b}$ (solid line).
		Figures from Daylan et al.~\cite{Daylan:2014rsa}.
	}
	\label{fig:gce}
\end{figure}

\textbf{Dwarf spheroidal galaxies.}
The Milky Way's satellite dwarf galaxies occupy the opposite end of the signal-versus-background trade-off.
Dwarf spheroidals are among the most dark matter--dominated objects known: their shallow gravitational potential wells allowed early supernova feedback to expel most of the baryonic gas, halting star formation and leaving behind systems with mass-to-light ratios ranging from tens of $M_\odot/L_\odot$ for classical dwarfs like Sculptor to thousands for ultrafaint dwarfs such as Segue~1 \cite{Cirelli:2024ssz}.
The scarcity of stars, gas, and other active astrophysical objects means that their gamma-ray backgrounds are negligible \cite{Cirelli:2024ssz}.
The expected dark matter flux from any individual dwarf is far smaller than from the Galactic Center --- roughly three orders of magnitude lower for a typical system \cite{Hooper:2024} --- but the absence of competing emission makes them the cleanest targets available.

Eight ``classical'' dwarf spheroidals (Sculptor, Fornax, Leo I and II, Ursa Minor, Draco, Carina, and Sextans) have been known since before the 1990s \cite{Cirelli:2024ssz}.
The advent of digital surveys --- SDSS, DES, and Gaia --- has brought the total count to nearly fifty, with many of the recent discoveries classified as ultrafaint dwarfs \cite{Cirelli:2024ssz}.
This growing census is expected to continue with the Rubin Observatory.

By combining observations of dozens of these systems in a stacked analysis, the Fermi-LAT collaboration has derived some of the tightest constraints on the annihilation cross-section to date \cite{Fermi-LAT:2015att,Hooper:2024}.
The analysis of 28 dwarf galaxies excludes the thermal relic cross-section ($\langle \sigma v_\text{rel} \rangle \simeq 2.2 \times 10^{-26}$ cm$^3$/s) for dark matter masses below $\sim 100$ GeV in the $b\bar{b}$ channel.
The dominant systematic uncertainty lies in the $J$-factors themselves, which must be inferred from the kinematics of the dwarfs' constituent stars.
For ultrafaint systems, the number of kinematic tracers can be extremely small, and removing or adding a single star from the sample can shift the reconstructed $J$-factor by orders of magnitude \cite{Cirelli:2024ssz}.
Additional systematics arise from the cusp--core ambiguity in the inner density profile and from the choice of integration angle $\theta_\star$ \cite{Cirelli:2024ssz}.
Despite these challenges, dwarfs remain the gold standard for model-independent constraints.
Chapters~\ref{ch:5} and~\ref{ch:6} of this thesis revisit the satellite galaxy programme with improved statistical methodology.

\textbf{Dark subhalos.} $\Lambda$CDM simulations predict far more gravitationally bound subhalos within the Milky Way's halo than the $\sim 50$ known satellite galaxies can account for \cite{Cirelli:2024ssz}.
Most of these subhalos are expected to be entirely dark --- devoid of stars and gas --- yet sufficiently massive to produce a detectable gamma-ray flux if the annihilation cross-section is near the thermal relic value \cite{Cirelli:2024ssz}.
Such dark subhalos would appear in the Fermi-LAT data as unassociated point sources with no counterpart at other wavelengths.
Their positions are unknown a priori, turning the search into a classification problem: distinguishing rare dark matter subhalos from the much larger population of unidentified astrophysical sources (primarily pulsars and faint blazars).
This challenge is addressed in Chapter~\ref{ch:5}, where machine learning methods are applied to classify unassociated Fermi-LAT sources in the presence of dataset shift between simulations and real data.

\textbf{The unresolved gamma-ray background.}
The cumulative, nearly isotropic gamma-ray flux that remains after subtracting resolved point sources and Galactic diffuse emission constitutes the unresolved gamma-ray background (UGRB) \cite{DGRB-review}.
This diffuse glow receives guaranteed contributions from unresolved blazars (BL~Lac objects and flat-spectrum radio quasars), misaligned active galactic nuclei, star-forming galaxies, and millisecond pulsars \cite{DGRB-review}.
Unresolved BL~Lac objects alone can account for the majority of the UGRB above $\sim 100$ GeV, while at lower energies other populations are required \cite{DGRB-review}.
A dark matter component --- originating from annihilation or decay in all halos and subhalos across cosmic history --- may also be present, but it is expected to be subdominant in the total intensity spectrum \cite{DGRB-review,Cirelli:2024ssz}.

Because the dark matter signal is buried beneath dominant astrophysical foregrounds, its extraction requires methods that go beyond the energy spectrum alone.
The angular anisotropies of the UGRB carry additional information: different source classes imprint different patterns of spatial fluctuations, reflecting their distinct clustering properties, redshift distributions, and luminosity functions \cite{DGRB-review}.
Cross-correlating these anisotropies with independent gravitational tracers of the large-scale structure --- galaxy catalogs, cosmic shear from weak lensing surveys, or $21\,$cm intensity mapping of neutral hydrogen --- can isolate the component that traces the dark matter distribution \cite{DGRB-review,Pinetti:2021jjs,Cirelli:2024ssz}.
The approach is particularly powerful because dark matter annihilation scales as $\rho^2$ and has a redshift evolution distinct from that of astrophysical gamma-ray emitters, providing a spectral and tomographic handle for separation \cite{Pinetti:2021jjs}.
This cross-correlation technique forms the basis of the analysis in Part~IV, where the sensitivity of next-generation instruments to the dark matter cross-correlation signal is forecast.

The complementary strengths of these targets --- the Galactic Center's raw flux, the dwarfs' clean backgrounds, the dark subhalos' model-independence, and the UGRB's cosmological reach --- motivate the multi-pronged programme developed in Parts II through IV of this thesis.
Each target exposes different aspects of the dark matter parameter space and is sensitive to different systematic uncertainties.
Progress requires pursuing all of them simultaneously, using the formalism established in this section as the common foundation.

\subsection{Multi-Messenger Context}
\label{sec:multi_messenger}

Dark matter annihilation and decay produce all stable Standard Model particles, not only gamma-rays.
Each additional messenger opens a complementary window onto the dark sector, but none offers the same combination of directional, spectral, and morphological information that makes gamma-rays the preferred channel for the analyses in this thesis \cite{Cirelli:2024ssz}.

\textbf{Neutrinos.}
Like photons, neutrinos are electrically neutral and propagate along straight lines, preserving the directional information that charged cosmic rays lose \cite{Cirelli:2024ssz,Hooper:2024}.
They can also escape environments that are completely opaque to photons.
The clearest example is the solar interior: dark matter particles from the Galactic halo occasionally scatter off nuclei in the Sun, lose enough kinetic energy to become gravitationally bound, and gradually sink to the core \cite{Cirelli:2024ssz}.
There they accumulate and annihilate, producing neutrinos that --- after oscillations and interactions within the dense solar medium --- are the only annihilation products that reach the surface \cite{Cirelli:2024ssz}.
A detection of high-energy neutrinos from the center of the Sun would therefore probe the dark matter scattering cross-section on nuclei, the same quantity measured by underground direct detection experiments \cite{Cirelli:2024ssz}.
Searches by IceCube, ANTARES, Super-Kamiokande, and the future KM3NeT and Hyper-Kamiokande have so far found no signal, but the resulting bounds on the spin-dependent scattering cross-section are approximately one order of magnitude stronger than those from direct detection experiments \cite{IceCube:2016dgk,Cirelli:2024ssz}.
For spin-independent scattering, the sensitivity is reversed: the neutrino constraints are about two orders of magnitude weaker than current direct detection limits \cite{Cirelli:2024ssz}.
In galactic searches --- from the Galactic Center, halo, and dwarf galaxies --- neutrino constraints are generally weaker than the corresponding gamma-ray bounds, but become competitive for dark matter masses above several TeV, because the neutrino interaction cross-section grows with energy and partly compensates for the reduced flux \cite{Cirelli:2024ssz}.

\textbf{Charged cosmic rays.}
Positrons, antiprotons, and heavier anti-nuclei are produced abundantly in dark matter annihilation, and the astrophysical backgrounds for antimatter are far smaller than for the corresponding matter species --- which is why dark matter searches concentrate on antiparticles rather than on the more abundant protons and electrons \cite{Cirelli:2024ssz}.
The fundamental limitation of all charged messengers is that galactic magnetic fields randomize their arrival directions, erasing the spatial information that makes gamma-ray searches so powerful and leaving only the energy spectrum as a diagnostic \cite{Cirelli:2024ssz,Pinetti:2021jjs}.

Positrons lose energy rapidly through synchrotron radiation and inverse Compton scattering as they propagate, so that the detected flux is dominated by local sources: roughly 90\% of 1 TeV positrons reaching Earth originate from within 1 kpc \cite{Cirelli:2024ssz}.
The steep rise in the cosmic-ray positron fraction above 10 GeV, first measured by PAMELA and confirmed with high precision by AMS-02, initially attracted considerable attention as a potential dark matter signal \cite{PAMELA:2008gwm,Cirelli:2024ssz}.
Fitting the excess with annihilating dark matter, however, requires a multi-TeV particle with a cross-section of order $\langle \sigma v_\text{rel} \rangle \sim 10^{-23}$ cm$^3$/s --- three orders of magnitude above the thermal relic value --- and purely leptophilic couplings to avoid overproducing antiprotons \cite{Cirelli:2024ssz}.
These requirements are in severe tension with gamma-ray constraints from dwarf galaxies and with CMB bounds on energy injection during recombination \cite{Cirelli:2024ssz}.
Today, the positron excess is widely attributed to nearby pulsars: the intense magnetic fields of rapidly spinning neutron stars strip electrons from the surface, producing $e^\pm$ pairs whose energy budget and spectral shape match the observed data without the tensions that plague the dark matter interpretation \cite{Cirelli:2024ssz,Hooper:2024}.

Antiprotons undergo negligible energy losses and thus preserve their injection spectrum, probing a much larger galactic volume than positrons \cite{Cirelli:2024ssz}.
AMS-02 data have revealed a potential excess of antiprotons at kinetic energies of 10--20 GeV, which can be fit by an $\sim$80 GeV dark matter particle annihilating into $b\bar{b}$ with a roughly thermal cross-section \cite{Cuoco:2016eej,Cirelli:2024ssz} --- properties tantalizingly close to those invoked to explain the Galactic Center gamma-ray excess (Chapter~\ref{ch:4}).
However, predictions of the low-energy antiproton spectrum are sensitive to the propagation model, the treatment of solar modulation, and the antiproton production cross-sections.
When these systematic uncertainties are properly accounted for, several comprehensive analyses find the data to be consistent with a purely secondary astrophysical origin, and the global significance of the excess drops below $\sim 2\sigma$ \cite{Cirelli:2024ssz}.

\textbf{Anti-nuclei.}
A qualitatively different and potentially decisive measurement would be the detection of anti-deuterons or anti-helium at low kinetic energies.
Heavier anti-nuclei are produced when antinucleons from a dark matter annihilation event coalesce --- that is, when an antiproton and an antineutron (or more antinucleons, for anti-helium) are emitted with momenta sufficiently close to merge into a bound state \cite{Cirelli:2024ssz}.
The yield is heavily suppressed, roughly a factor of $\mathcal{O}(10^{-5})$ per additional antinucleon, but the signal benefits from an almost vanishing astrophysical background \cite{Cirelli:2024ssz}.
Standard secondary anti-deuterons, produced by spallation of high-energy cosmic rays on interstellar gas, peak at a few GeV per nucleon, while dark matter annihilation --- occurring essentially at rest --- produces anti-nuclei that peak below 1 GeV per nucleon, a regime where the background is practically zero \cite{Cirelli:2024ssz,Hooper:2024}.
Current bounds from the BESS balloon experiment are too weak to constrain dark matter \cite{Cirelli:2024ssz}, but the projected sensitivities of AMS-02 and the upcoming GAPS experiment --- a balloon-borne detector that identifies anti-deuterons by capturing them into exotic atoms whose annihilation emits characteristic X-rays and pions --- may reach the thermal relic cross-section for 100 GeV dark matter under favorable conditions \cite{Cirelli:2024ssz}.
AMS-02 has also tentatively reported $\sim$9 anti-helium candidate events \cite{Cirelli:2024ssz}.
If confirmed, this observation would be difficult to explain with either standard astrophysical production or current dark matter models, both of which predict fluxes many orders of magnitude below the detector sensitivity \cite{Cirelli:2024ssz}.

This thesis focuses exclusively on \textbf{gamma-rays} --- the messenger that provides both spatial and spectral information while relying on the fewest astrophysical assumptions.
The multi-messenger context is noted where relevant, but the analyses of Parts II--V are rooted entirely in the gamma-ray formalism developed in this section.

\subsection{The Status of the Field}
\label{sec:status}

Despite more than a decade of increasingly sensitive searches across all three detection strategies, no unambiguous signal of dark matter has been found \cite{Cirelli:2024ssz}.
The current status is best summarized as a set of progressively tightening upper bounds on $\langle \sigma v_\text{rel} \rangle$ and $\tau$, punctuated by anomalies whose dark matter interpretations have been challenged by plausible astrophysical alternatives.

\textbf{Gamma-ray constraints.}
The tightest indirect detection limits come from the Fermi-LAT stacked analysis of Milky Way dwarf spheroidal galaxies.
By combining observations of nearly thirty dSphs --- systems with negligible astrophysical backgrounds and dark matter mass-to-light ratios reaching into the thousands --- the Fermi-LAT collaboration excludes the thermal relic cross-section for dark matter masses below approximately 100 GeV in the $b\bar{b}$ channel \cite{Fermi-LAT:2015att,Hooper:2024}.
For a generic annihilation channel, marginalizing over branching ratios yields a model-independent lower limit of $m_\chi \gtrsim 20$ GeV \cite{Cirelli:2024ssz}.
Constraints on the dark matter lifetime from gamma-ray observations reach $\tau \gtrsim 10^{27}$--$10^{28}$ seconds --- some ten orders of magnitude longer than the age of the universe \cite{Cirelli:2024ssz}.
The central target for these searches --- the thermal relic cross-section --- is illustrated in Figure~\ref{fig:thermal_relic}, which shows how the freeze-out abundance varies with the dark matter mass and annihilation cross-section.

\begin{figure}[t]
	\centering
	\includegraphics[width=0.47\textwidth]{figures/Omegafreeze-out.pdf}
	\caption{Dark matter freeze-out abundance $\Omega_\text{DM} h^2$ as a function of DM mass and annihilation cross-section $\langle \sigma v \rangle$ for Majorana dark matter.
		The measured cosmological abundance $\Omega_\text{DM} h^2 = 0.120 \pm 0.001$ is reproduced within the orange band at $3\sigma$.
		For smaller (larger) annihilation cross-sections --- i.e.\ in the lower (upper) part of the plot --- the freeze-out happens too early (too late) and therefore the DM relic abundance is too high (too low).
		This cross-section is the target for indirect detection searches.
		Figure from Cirelli et al.~\cite{Cirelli:2024ssz}.
	}
	\label{fig:thermal_relic}
\end{figure}

\textbf{Direct detection and collider constraints.}
The tightening of parameter space is not confined to indirect searches.
The current generation of dual-phase liquid xenon detectors --- LZ, XENONnT, and PandaX-4T --- has pushed the spin-independent WIMP--nucleon cross-section below $10^{-46}$ cm$^2$ for masses near the electroweak scale \cite{Cirelli:2024ssz,Hooper:2024}.
At the LHC, mono-jet and mono-photon searches have found no evidence for new physics, excluding strongly interacting superpartners (squarks, gluinos) up to masses of 1--2 TeV and charged superpartners up to several hundred GeV \cite{Cirelli:2024ssz,Hooper:2024}.
The theoretical basis of the most popular WIMP candidate --- the supersymmetric neutralino --- has been substantially weakened by the absence of supersymmetry at collider energies \cite{Cirelli:2024ssz}.

\textbf{Unresolved anomalies.}
Against this backdrop of null results, several persistent anomalies keep the possibility of discovery alive.
The Galactic Center gamma-ray excess --- a signal first reported in 2009 and consistent with a $\sim$40--70 GeV dark matter particle annihilating into $b\bar{b}$ at the thermal cross-section --- remains the most debated \cite{Daylan:2014rsa,Cirelli:2024ssz,Hooper:2024}.
Its leading astrophysical competitor, an unresolved population of millisecond pulsars in the Galactic bulge, has been alternately supported and challenged by photon-count statistics, morphological analyses, and luminosity function arguments \cite{Cirelli:2024ssz}.
Chapter~\ref{ch:4} of this thesis contributes directly to this debate.
In charged cosmic rays, the AMS-02 antiproton data show a tentative excess at kinetic energies around 10--20 GeV that could be fit by an 80 GeV dark matter particle annihilating into $b\bar{b}$ --- parameters that overlap suggestively with those of the Galactic Center excess \cite{Cirelli:2024ssz}.
However, once the substantial systematic uncertainties in the cosmic-ray propagation models and antiproton production cross-sections are accounted for, the significance of this excess drops below $2\sigma$, and the data remain consistent with a purely secondary astrophysical origin \cite{Cirelli:2024ssz}.
The positron excess measured by PAMELA and AMS-02 above 10 GeV, discussed in Section~\ref{sec:multi_messenger}, is now widely attributed to nearby pulsars rather than dark matter \cite{Cirelli:2024ssz,Hooper:2024}.

\textbf{The methodological case for this thesis.}
Taken together, these results define the context in which this thesis operates.
The viable parameter space for the simplest thermal WIMP models has narrowed substantially --- what Cirelli et al.
describe as a field ``at a turning point'' \cite{Cirelli:2024ssz}.
The standard approach to indirect detection --- searching for a bright, spectrally distinctive excess from a single target above the astrophysical background --- has not delivered a discovery.
Yet the thermal relic benchmark is not the prediction of a specific model but a cosmological target: if dark matter is a thermal relic, the annihilation cross-section must have the value that produces the observed abundance, regardless of the underlying theory.
Current exclusions cover only a portion of the mass range and a subset of annihilation channels; heavier candidates, leptonic final states, and velocity-dependent cross-sections remain largely unexplored.
Extracting these signals demands methods that go beyond simple signal-over-background thresholding.

The analyses of this thesis pursue that transition.
The Galactic Center excess demands careful modelling of the interstellar emission and the unresolved point-source population --- distinguishing dark matter annihilation from millisecond pulsars requires statistical tools, not a single spectral fit (Part~II).
Individual dark matter subhalos among the Fermi-LAT unassociated sources cannot be identified by classification alone; a fully probabilistic model that accounts for both covariate and prior probability shifts between associated and unassociated sources is needed to derive statistically rigorous constraints (Part~II).
Below the point-source detection threshold, deep learning applied to gamma-ray sky maps can recover the source-count distribution of faint populations down to fluxes a factor of 50 below the Fermi-LAT sensitivity limit, while probabilistic catalogues extend spatial source information into the sub-threshold regime (Part~III).
At the largest scales, the dark matter contribution to the unresolved gamma-ray background can be isolated through its cross-correlation with independent tracers of the large-scale structure (Part~IV).

The mathematical framework developed in this section --- the flux factorization, the $J$-factor and $D$-factor formalism, the density profiles and their uncertainties --- underpins every one of these analyses.

